%Version 3.1 December 2024
% See section 11 of the User Manual for version history
%
%%%%%%%%%%%%%%%%%%%%%%%%%%%%%%%%%%%%%%%%%%%%%%%%%%%%%%%%%%%%%%%%%%%%%%
%%                                                                 %%
%% Please do not use \input{...} to include other tex files.       %%
%% Submit your LaTeX manuscript as one .tex document.              %%
%%                                                                 %%
%% All additional figures and files should be attached             %%
%% separately and not embedded in the \TeX\ document itself.       %%
%%                                                                 %%
%%%%%%%%%%%%%%%%%%%%%%%%%%%%%%%%%%%%%%%%%%%%%%%%%%%%%%%%%%%%%%%%%%%%%

%%\documentclass[referee,sn-basic]{sn-jnl}% referee option is meant for double line spacing

%%=======================================================%%
%% to print line numbers in the margin use lineno option %%
%%=======================================================%%

%%\documentclass[lineno,pdflatex,sn-basic]{sn-jnl}% Basic Springer Nature Reference Style/Chemistry Reference Style

%%=========================================================================================%%
%% the documentclass is set to pdflatex as default. You can delete it if not appropriate.  %%
%%=========================================================================================%%

%%\documentclass[sn-basic]{sn-jnl}% Basic Springer Nature Reference Style/Chemistry Reference Style

%%Note: the following reference styles support Namedate and Numbered referencing. By default the style follows the most common style. To switch between the options you can add or remove “Numbered” in the optional parenthesis. 
%%The option is available for: sn-basic.bst, sn-chicago.bst%  
 
%%\documentclass[pdflatex,sn-nature]{sn-jnl}% Style for submissions to Nature Portfolio journals
%%\documentclass[pdflatex,sn-basic]{sn-jnl}% Basic Springer Nature Reference Style/Chemistry Reference Style
\documentclass[pdflatex,sn-mathphys-num]{sn-jnl}% Math and Physical Sciences Numbered Reference Style
%%\documentclass[pdflatex,sn-mathphys-ay]{sn-jnl}% Math and Physical Sciences Author Year Reference Style
%%\documentclass[pdflatex,sn-aps]{sn-jnl}% American Physical Society (APS) Reference Style
%%\documentclass[pdflatex,sn-vancouver-num]{sn-jnl}% Vancouver Numbered Reference Style
%%\documentclass[pdflatex,sn-vancouver-ay]{sn-jnl}% Vancouver Author Year Reference Style
%%\documentclass[pdflatex,sn-apa]{sn-jnl}% APA Reference Style
%%\documentclass[pdflatex,sn-chicago]{sn-jnl}% Chicago-based Humanities Reference Style

%%%% Standard Packages
%%<additional latex packages if required can be included here>

\usepackage{graphicx}%
\usepackage{multirow}%
\usepackage{amsmath,amssymb,amsfonts}%
\usepackage{amsthm}%
\usepackage{mathrsfs}%
\usepackage[title]{appendix}%
\usepackage{xcolor}%
\usepackage{textcomp}%
\usepackage{manyfoot}%
\usepackage{booktabs}%
\usepackage{algorithm}%
\usepackage{algorithmicx}%
\usepackage{algpseudocode}%
\usepackage{listings}%
%%%%

%%%%%=============================================================================%%%%
%%%%  Remarks: This template is provided to aid authors with the preparation
%%%%  of original research articles intended for submission to journals published 
%%%%  by Springer Nature. The guidance has been prepared in partnership with 
%%%%  production teams to conform to Springer Nature technical requirements. 
%%%%  Editorial and presentation requirements differ among journal portfolios and 
%%%%  research disciplines. You may find sections in this template are irrelevant 
%%%%  to your work and are empowered to omit any such section if allowed by the 
%%%%  journal you intend to submit to. The submission guidelines and policies 
%%%%  of the journal take precedence. A detailed User Manual is available in the 
%%%%  template package for technical guidance.
%%%%%=============================================================================%%%%

%% as per the requirement new theorem styles can be included as shown below
\theoremstyle{thmstyleone}%
\newtheorem{theorem}{Theorem}%  meant for continuous numbers
%%\newtheorem{theorem}{Theorem}[section]% meant for sectionwise numbers
%% optional argument [theorem] produces theorem numbering sequence instead of independent numbers for Proposition
\newtheorem{proposition}[theorem]{Proposition}% 
%%\newtheorem{proposition}{Proposition}% to get separate numbers for theorem and proposition etc.

\theoremstyle{thmstyletwo}%
\newtheorem{example}{Example}%
\newtheorem{remark}{Remark}%

\theoremstyle{thmstylethree}%
\newtheorem{definition}{Definition}%

\raggedbottom
%%\unnumbered% uncomment this for unnumbered level heads

\begin{document}

\title[Article Title]{Article Title}

%%=============================================================%%
%% GivenName	-> \fnm{Joergen W.}
%% Particle	-> \spfx{van der} -> surname prefix
%% FamilyName	-> \sur{Ploeg}
%% Suffix	-> \sfx{IV}
%% \author*[1,2]{\fnm{Joergen W.} \spfx{van der} \sur{Ploeg} 
%%  \sfx{IV}}\email{iauthor@gmail.com}
%%=============================================================%%

\author*[1]{\fnm{Lalith} \sur{Gunnu}}\email{lalithgunnu2202@gmail.com}

\author[2]{\fnm{Kaustab} \sur{Karmakar}}

\affil*[1]{\orgdiv{Department of Electronics and Instrumentation Engineering}, \orgname{National Institute of Technology Silchar}, \orgaddress{\city{Silchar}, \state{Assam}, \country{India}}}

\affil[2]{\orgdiv{Department of Computer Science and Engineering}, \orgname{National Institute of Technology Silchar}, \orgaddress{\city{Silchar}, \state{Assam}, \country{India}}}

%%==================================%%
%% Sample for unstructured abstract %%
%%==================================%%

\abstract{The abstract serves both as a general introduction to the topic and as a brief, non-technical summary of the main results and their implications. Authors are advised to check the author instructions for the journal they are submitting to for word limits and if structural elements like subheadings, citations, or equations are permitted.}

%%================================%%
%% Sample for structured abstract %%
%%================================%%

% \abstract{\textbf{Purpose:} The abstract serves both as a general introduction to the topic and as a brief, non-technical summary of the main results and their implications. The abstract must not include subheadings (unless expressly permitted in the journal's Instructions to Authors), equations or citations. As a guide the abstract should not exceed 200 words. Most journals do not set a hard limit however authors are advised to check the author instructions for the journal they are submitting to.
% 
% \textbf{Methods:} The abstract serves both as a general introduction to the topic and as a brief, non-technical summary of the main results and their implications. The abstract must not include subheadings (unless expressly permitted in the journal's Instructions to Authors), equations or citations. As a guide the abstract should not exceed 200 words. Most journals do not set a hard limit however authors are advised to check the author instructions for the journal they are submitting to.
% 
% \textbf{Results:} The abstract serves both as a general introduction to the topic and as a brief, non-technical summary of the main results and their implications. The abstract must not include subheadings (unless expressly permitted in the journal's Instructions to Authors), equations or citations. As a guide the abstract should not exceed 200 words. Most journals do not set a hard limit however authors are advised to check the author instructions for the journal they are submitting to.
% 
% \textbf{Conclusion:} The abstract serves both as a general introduction to the topic and as a brief, non-technical summary of the main results and their implications. The abstract must not include subheadings (unless expressly permitted in the journal's Instructions to Authors), equations or citations. As a guide the abstract should not exceed 200 words. Most journals do not set a hard limit however authors are advised to check the author instructions for the journal they are submitting to.}

\keywords{keyword1, Keyword2, Keyword3, Keyword4}

%%\pacs[JEL Classification]{D8, H51}

%%\pacs[MSC Classification]{35A01, 65L10, 65L12, 65L20, 65L70}

\maketitle

\section{Introduction}\label{sec1}

The Introduction section, of referenced text \cite{bib1} expands on the background of the work (some overlap with the Abstract is acceptable). The introduction should not include subheadings.

Springer Nature does not impose a strict layout as standard however authors are advised to check the individual requirements for the journal they are planning to submit to as there may be journal-level preferences. When preparing your text please also be aware that some stylistic choices are not supported in full text XML (publication version), including coloured font. These will not be replicated in the typeset article if it is accepted. 

\section{Results}\label{sec2}

Sample body text. Sample body text. Sample body text. Sample body text. Sample body text. Sample body text. Sample body text. Sample body text.

```latex
\section{Research Questions}\label{sec3}

\subsection{Construction of Shared Multimodal Embedding Space}\label{subsec3.1}

The construction of a shared multimodal embedding space is a fundamental requirement for enabling cross-modal representation learning, where information from different modalities is mapped into a common feature space. The primary objective is to establish semantic correspondence between modality-specific representations, allowing features originating from different modalities to be compared directly. Across CLIP, AudioCLIP, and CLAP, this is achieved through modality-specific neural network encoders, learnable linear projection layers, and contrastive learning objectives. Although these approaches differ in the number and nature of the modalities involved, they share a common principle: representations of semantically corresponding inputs are brought closer in the embedding space, while representations of non-corresponding inputs are pushed apart.

\subsubsection{Modality-Specific Encoders and Feature Projection}

The initial stage in constructing a shared embedding space involves extracting meaningful representations from each input modality using dedicated neural network encoders. In CLIP, Radford et al. employ a dual-encoder architecture consisting of a vision encoder and a Transformer-based text encoder. The visual encoder utilizes either a modified ResNet with attention pooling or a Vision Transformer (ViT), while the text encoder processes byte-pair encoded textual sequences using a Transformer architecture. The resulting feature representations are subsequently mapped through learnable linear projection layers into a common 1024-dimensional embedding space. This projection enables image and text representations to be expressed in a unified vector space despite originating from different encoder architectures.

Extending this framework to three modalities, AudioCLIP incorporates a dedicated audio encoder alongside the pre-existing CLIP image and text encoders. The audio branch employs an ESResNeXt-50 architecture with trainable B-spline wavelet-based time-frequency transformations to extract acoustic representations. Its classification layer is replaced with a linear projection layer that maps the extracted audio features into the same 1024-dimensional embedding space used by the image and text encoders. Consequently, visual, textual, and acoustic representations become directly comparable within a common multimodal space.

Similarly, CLAP adopts a dual-encoder architecture specifically designed for audio-text representation learning. It combines a CNN14-based audio encoder, which processes log-Mel spectrogram representations, with a BERT-base text encoder that extracts semantic representations from natural language descriptions using the [CLS] token embedding. The resulting audio and textual features are independently transformed using learnable projection matrices, mapping both representations into a shared 1024-dimensional embedding space. Unlike AudioCLIP, which extends an existing image-text alignment framework, CLAP directly establishes correspondence between acoustic signals and natural language descriptions.

\subsubsection{Contrastive Learning for Cross-Modal Alignment}

Following feature extraction and projection, contrastive learning is employed to establish semantic alignment between the representations of different modalities. In CLIP, this is achieved by constructing an $N \times N$ similarity matrix for a batch containing $N$ matched image-text pairs. Each element represents the cosine similarity between an image embedding and a text embedding. A symmetric cross-entropy objective is then applied across both image-to-text and text-to-image directions. The objective maximizes the similarity of corresponding image-text pairs, represented by the diagonal elements of the similarity matrix, while minimizing the similarity of non-matching pairs represented by the off-diagonal elements. A learnable temperature parameter is used to scale the similarity logits, controlling the sharpness of the resulting probability distribution.

AudioCLIP extends this contrastive alignment mechanism to a tri-modal setting by introducing additional cross-modal objectives. In addition to the original image-text alignment objective, it incorporates text-audio and image-audio contrastive learning objectives. This enables the model to establish correspondence across all three modalities rather than restricting alignment to a single modality pair. AudioSet provides the multimodal training associations by connecting audio recordings with corresponding video frames and textual class labels. Its training strategy involves standalone pre-training of the audio encoder, cooperative pre-training using the pre-trained image and text encoders as fixed reference representations, and subsequent joint optimization of the three modality-specific branches. This progressive training procedure facilitates the integration of acoustic representations into the existing visual-textual embedding space.

In CLAP, cross-modal alignment is established directly between audio and text embeddings using a pairwise similarity matrix. For a batch of $N$ matched audio-text pairs, the similarity matrix is computed using the projected audio and text representations, with a learnable temperature parameter scaling the similarity scores. Similar to CLIP, a symmetric cross-entropy loss is applied in both audio-to-text and text-to-audio directions. This objective encourages corresponding audio recordings and textual descriptions to occupy nearby positions in the embedding space while distinguishing them from unrelated pairs. An important distinction is the use of natural language supervision derived from audio-captioning datasets, including AudioCaps, ClothoV2, FSD50K, and MACS. This provides richer semantic supervision than approaches relying exclusively on predefined class labels.

\subsubsection{Comparative Analysis of Shared Embedding Construction}

The three approaches demonstrate a progressive extension of contrastive multimodal representation learning. CLIP establishes the foundational image-text alignment framework through modality-specific encoders, linear feature projections, and symmetric contrastive optimization. AudioCLIP extends this framework by incorporating an additional acoustic encoder and introducing cross-modal objectives that connect audio with both visual and textual representations. In contrast, CLAP adopts a dedicated audio-text architecture that focuses on aligning acoustic information directly with natural language descriptions.

Despite these architectural differences, all three approaches rely on a common embedding dimensionality, learnable projection layers, and similarity-based contrastive objectives to establish cross-modal correspondence. However, their training supervision and alignment strategies differ according to the modalities being integrated. CLIP relies on large-scale image-text associations, AudioCLIP utilizes audio-visual-textual associations to connect three modalities, whereas CLAP leverages audio-caption pairs to establish direct semantic correspondence between sound and language. These approaches collectively demonstrate how contrastive learning provides a scalable mechanism for constructing shared multimodal embedding spaces, forming the foundation for subsequent cross-modal retrieval, recognition, and multimodal understanding tasks.
```latex
\section{Architectural Frameworks and Training Strategies}\label{sec5}

This section examines the architectural configurations and training strategies adopted by CLIP, AudioCLIP, and CLAP. Although all three frameworks employ contrastive pre-training to establish shared multimodal embedding spaces, they differ in their encoder architectures, feature extraction mechanisms, modality integration strategies, and training procedures.

\subsection{CLIP: Visual--Textual Representation Learning}

\subsubsection{Encoder Architectures}

CLIP employs two independent neural network encoders to extract visual and textual representations. The image encoder evaluates two primary architectural families: Modified ResNet and Vision Transformer (ViT).

The Modified ResNet family includes ResNet-50, ResNet-101, and scaled variants such as RN50x4, RN50x16, and RN50x64. These architectures incorporate ResNet-D modifications and anti-aliased blur pooling to improve feature extraction. Additionally, the conventional global average pooling operation is replaced with a multi-head attention pooling mechanism based on query, key, and value (QKV) representations.

The Vision Transformer family comprises ViT-B/32, ViT-B/16, ViT-L/14, and ViT-L/14@336px. These architectures process images as sequences of patches using transformer-based self-attention. An additional layer normalization operation is applied before the transformer to facilitate stable representation learning.

For textual representation, CLIP employs a Transformer-based language encoder. Its base configuration consists of 12 transformer layers, a hidden representation width of 512, and 8 attention heads. Input text is converted into lowercase byte-pair encoding (BPE) tokens using a vocabulary of 49,408 tokens, with a maximum sequence length of 76 tokens. The representation corresponding to the end-of-sequence ([EOS]) token is layer-normalized and subsequently projected into the shared multimodal embedding space.

Both image and text encoders employ learnable linear projection layers to transform their respective feature representations into a common embedding dimension, enabling direct cross-modal similarity computation.

\subsubsection{Contrastive Pre-training Strategy}

CLIP is trained using a contrastive learning objective based on symmetric cross-entropy, commonly associated with InfoNCE and multi-class N-pair learning. Given a minibatch of $N$ image--text pairs, the model computes an $N \times N$ similarity matrix and identifies the corresponding image--text pairs among all possible combinations.

The diagonal elements of the similarity matrix represent positive image--text associations, whereas the off-diagonal elements represent negative associations. A learnable temperature parameter controls the scaling of similarity logits. The temperature is initialized to $0.07$, and the resulting logit scaling factor is constrained to a maximum value of $100$.

The training objective minimizes the average of the image-to-text and text-to-image cross-entropy losses, encouraging the model to establish semantically aligned visual and textual representations within the shared embedding space.

\subsection{AudioCLIP: Tri-Modal Representation Learning}

\subsubsection{Encoder Architectures}

AudioCLIP extends the original CLIP architecture by incorporating an additional audio processing branch, thereby enabling joint representation learning across image, text, and audio modalities.

The image and text encoders are initialized using the pre-trained CLIP architecture. Specifically, the image branch employs a ResNet-50 backbone equipped with QKV attention pooling, while the textual branch retains the pre-trained Transformer-based language encoder.

For audio representation, AudioCLIP integrates ESResNeXt-50, a convolutional neural network architecture containing approximately 30 million parameters. Unlike conventional audio classification networks, ESResNeXt incorporates a trainable time-frequency transformation based on complex frequency B-spline wavelets. This transformation enables the extraction of discriminative acoustic representations from input audio signals.

The original classification layer of ESResNeXt is replaced with a learnable linear projection layer that maps audio representations into the same 1024-dimensional embedding space used by the image and text branches. This architectural modification enables direct similarity computation across all three modalities.

\subsubsection{Multi-Stage Training Strategy}

AudioCLIP adopts a multi-stage training procedure to integrate the audio encoder with the pre-trained visual--textual representation space.

\begin{enumerate}
    \item \textbf{Standalone Audio Pre-training:}
    The ESResNeXt audio encoder is initially pre-trained on the AudioSet dataset to learn general-purpose acoustic representations. This stage establishes the audio encoder's ability to extract meaningful features from diverse sound categories.

    \item \textbf{Cooperative Pre-training:}
    The pre-trained audio encoder is integrated with the existing CLIP image and text branches. During this stage, the image and text encoders remain frozen and act as teacher networks, while the audio branch is optimized to align its representations with the established visual--textual embedding space. This procedure facilitates cross-modal knowledge transfer into the audio representation space.

    \item \textbf{Full Model Fine-tuning:}
    Following cooperative pre-training, all three modality-specific encoders are jointly optimized using audio clips, corresponding video frames, and associated class names from AudioSet. This stage enables the model to refine inter-modal relationships and establish a unified tri-modal representation space.
\end{enumerate}

The training objective combines three pairwise contrastive losses corresponding to image--text, text--audio, and image--audio associations. Consequently, AudioCLIP extends the original bidirectional alignment mechanism of CLIP into a tri-modal learning framework.

\subsection{CLAP: Language--Audio Representation Learning}

\subsubsection{Encoder Architectures}

CLAP employs separate audio and textual encoders to establish a shared language--audio representation space. Unlike AudioCLIP, which extends an existing visual--textual framework, CLAP is specifically designed to learn semantic correspondences between natural language descriptions and acoustic representations.

The audio encoder utilizes CNN14, a convolutional neural network architecture derived from the PANNs framework, containing approximately 80.8 million parameters. Input audio is represented using log-Mel spectrograms generated at a sampling rate of 44.1 kHz, with 64 Mel frequency bins, a hop size of 320 samples, and a window size of 1024 samples. The CNN14 architecture extracts 2048-dimensional audio feature representations.

For textual feature extraction, CLAP employs BERT-base uncased, consisting of approximately 110 million parameters. The contextual representation corresponding to the [CLS] token from the final transformer layer is extracted as the textual feature vector, producing a 768-dimensional representation. The input text is processed using a maximum character length of 100.

The audio and text representations are subsequently transformed into a shared 1024-dimensional embedding space through two independent learnable linear projection matrices, denoted by $L_a$ and $L_t$, respectively. This projection mechanism enables the computation of cross-modal similarity between acoustic and textual representations.

\subsubsection{Contrastive Pre-training Strategy}

CLAP adopts a symmetric contrastive cross-entropy objective to align corresponding audio and text representations. For a minibatch containing $N$ audio--text pairs, pairwise cosine similarities are computed to construct an $N \times N$ similarity matrix.

The similarity logits are scaled using a learnable temperature parameter, with an initial value reported as $0.007$ in the considered configuration. The scaling factor is constrained to maintain numerical stability during optimization.

Unlike classification-based approaches that require predefined acoustic categories, CLAP employs natural language supervision to learn semantic relationships directly from audio-caption pairs. Its contrastive objective minimizes both text-to-audio and audio-to-text cross-entropy losses, enabling bidirectional alignment between the two modalities.

\subsubsection{Training Configuration and Dataset Composition}

Both the audio and text encoders are jointly optimized during contrastive pre-training. The training configuration considered here utilizes approximately 128,010 audio--text pairs collected from four audio captioning datasets: FSD50K, ClothoV2, AudioCaps, and MACS.

The combination of these datasets provides diverse acoustic environments and natural language descriptions, allowing the model to learn semantic associations across different sound categories and descriptive expressions.

Through joint optimization of the two encoders, CLAP establishes a shared representation space that supports cross-modal retrieval and zero-shot audio classification without requiring task-specific classifier training.

\subsection{Comparative Architectural Analysis}

The three frameworks exhibit distinct architectural and training characteristics despite sharing a contrastive learning foundation. CLIP establishes visual--textual alignment through independently encoded image and language representations. AudioCLIP extends this mechanism by incorporating an audio encoder and introducing pairwise alignment across three modalities. In contrast, CLAP focuses exclusively on audio--language alignment through dedicated acoustic and textual encoders trained using natural language supervision.

Their architectural configurations and training strategies demonstrate different approaches to multimodal representation learning, particularly in terms of modality integration, encoder initialization, and the utilization of pre-trained representations.
```

\section{Mathematical Formulation of Contrastive Multimodal Models}\label{sec4}

This section presents the mathematical formulations of three representative contrastive multimodal learning frameworks: Contrastive Language--Image Pre-training (CLIP), AudioCLIP, and Contrastive Language--Audio Pre-training (CLAP). These frameworks learn shared embedding spaces by maximizing the similarity between corresponding multimodal representations while minimizing the similarity between non-corresponding samples. Their respective feature extraction mechanisms, projection operations, similarity computations, and optimization objectives are described below.

\subsection{CLIP: Contrastive Language--Image Pre-training}

\subsubsection{Feature Extraction and Multimodal Projection}

Given a minibatch of $N$ paired images and text descriptions, denoted by $\{(I_i,T_i)\}_{i=1}^{N}$, CLIP employs separate image and text encoders to extract their respective feature representations. The extracted features are expressed as

\begin{align}
I_f &= f_{\mathrm{image}}(I) \in \mathbb{R}^{N \times d_i}, \\
T_f &= f_{\mathrm{text}}(T) \in \mathbb{R}^{N \times d_t}.
\end{align}

The image and text representations are subsequently projected into a common multimodal embedding space of dimension $d_e$ using learnable linear projection matrices $W_i$ and $W_t$, respectively. The resulting embeddings are normalized to unit length:

\begin{align}
I_e &= \operatorname{normalize}(I_f W_i)
\in \mathbb{R}^{N \times d_e}, \\
T_e &= \operatorname{normalize}(T_f W_t)
\in \mathbb{R}^{N \times d_e}.
\end{align}

Here, $W_i \in \mathbb{R}^{d_i \times d_e}$ and $W_t \in \mathbb{R}^{d_t \times d_e}$ are learnable projection matrices. The normalization operation ensures that the resulting embeddings have unit $L_2$ norm, enabling cosine similarity to be computed through their dot product.

\subsubsection{Similarity Computation and Temperature Scaling}

CLIP computes pairwise similarities between all image and text embeddings within a minibatch, producing a similarity logit matrix $S \in \mathbb{R}^{N \times N}$. Each element represents the scaled cosine similarity between an image embedding and a text embedding:

\begin{equation}
S_{ij} = \frac{I_{e,i} \cdot T_{e,j}^{T}}{\tau}
= \exp(t)(I_{e,i} \cdot T_{e,j}^{T}),
\label{eq:clip_similarity}
\end{equation}

where $\tau$ denotes the learnable temperature parameter and $t=\log(1/\tau)$ is its logarithmic representation. The temperature is initialized to $0.07$, and the exponential logit scale is constrained to a maximum value of $100$ during optimization to maintain numerical stability.

\subsubsection{Contrastive Loss Function}

CLIP employs a symmetric contrastive learning objective based on cross-entropy loss. For a minibatch of $N$ image--text pairs, the diagonal elements of the similarity matrix represent the corresponding positive pairs, whereas the off-diagonal elements represent negative pairs.

The image-to-text contrastive loss is defined as

\begin{equation}
\mathcal{L}_{I \rightarrow T}
=
-\frac{1}{N}
\sum_{i=1}^{N}
\log
\frac{\exp(S_{ii})}
{\sum_{j=1}^{N}\exp(S_{ij})}.
\label{eq:clip_i2t}
\end{equation}

Similarly, the text-to-image contrastive loss is expressed as

\begin{equation}
\mathcal{L}_{T \rightarrow I}
=
-\frac{1}{N}
\sum_{j=1}^{N}
\log
\frac{\exp(S_{jj})}
{\sum_{i=1}^{N}\exp(S_{ij})}.
\label{eq:clip_t2i}
\end{equation}

The overall CLIP objective is obtained by averaging the two directional losses:

\begin{equation}
\mathcal{L}_{\mathrm{CLIP}}
=
\frac{1}{2}
\left(
\mathcal{L}_{I \rightarrow T}
+
\mathcal{L}_{T \rightarrow I}
\right).
\label{eq:clip_total}
\end{equation}

This symmetric objective encourages corresponding image and text representations to occupy nearby positions in the shared embedding space while separating unrelated pairs.

\subsection{AudioCLIP: Tri-Modal Contrastive Learning}

\subsubsection{Feature Representation and Audio Projection}

AudioCLIP extends the CLIP framework by incorporating an additional audio encoder, enabling joint representation learning across image, text, and audio modalities. While the image and text branches retain the original CLIP encoders, the audio branch employs an ESResNeXt architecture to extract acoustic representations.

For an input audio minibatch $A$, the audio encoder generates feature representations that are projected into the shared multimodal embedding space:

\begin{equation}
A_e =
\operatorname{normalize}
\left(f_{\mathrm{audio}}(A)W_a\right)
\in \mathbb{R}^{N \times d_e},
\label{eq:audioclip_audio}
\end{equation}

where $W_a$ denotes the learnable audio projection matrix. The resulting audio embeddings are normalized to maintain compatibility with the image and text representations in the common embedding space.

\subsubsection{Pairwise Similarity Matrices}

With three modality-specific embedding representations, AudioCLIP computes pairwise similarity matrices for all possible modality combinations. These matrices are defined as

\begin{align}
S^{(I,T)} &= I_e T_e^T \exp(t), \\
S^{(T,A)} &= T_e A_e^T \exp(t), \\
S^{(I,A)} &= I_e A_e^T \exp(t).
\label{eq:audioclip_similarity}
\end{align}

Each matrix captures the pairwise semantic correspondence between two modalities. The diagonal elements represent matching samples, while the remaining elements represent non-matching combinations within the minibatch.

\subsubsection{Tri-Modal Contrastive Objective}

AudioCLIP extends the symmetric contrastive objective of CLIP by incorporating the losses associated with all three modality pairs. For any two modalities $M_1$ and $M_2$, the corresponding symmetric contrastive loss is defined as

\begin{align}
\mathcal{L}_{(M_1,M_2)}
=
\frac{1}{2}
\Bigg[
&-\frac{1}{N}
\sum_{i=1}^{N}
\log
\frac{\exp(S_{ii}^{(M_1,M_2)})}
{\sum_{j=1}^{N}\exp(S_{ij}^{(M_1,M_2)})}
\nonumber\\
&-\frac{1}{N}
\sum_{j=1}^{N}
\log
\frac{\exp(S_{jj}^{(M_1,M_2)})}
{\sum_{i=1}^{N}\exp(S_{ij}^{(M_1,M_2)})}
\Bigg].
\label{eq:audioclip_pair_loss}
\end{align}

The overall tri-modal contrastive objective is then expressed as

\begin{equation}
\begin{aligned}
\mathcal{L}_{\mathrm{AudioCLIP}}
={}&
\mathcal{L}_{(I,T)}
+
\mathcal{L}_{(T,A)}
+
\mathcal{L}_{(I,A)}.
\end{aligned}
\label{eq:audioclip_total}
\end{equation}

The formulation enables the model to learn cross-modal correspondences not only between image and text, but also between audio and the other two modalities. During cooperative pre-training, the pretrained image and text branches act as teacher representations while the audio branch is optimized. In subsequent full-model training, the participating modality encoders can be jointly optimized.

\subsection{CLAP: Contrastive Language--Audio Pre-training}

\subsubsection{Audio and Text Feature Extraction}

CLAP learns a shared embedding space for audio and natural language representations through separate audio and text encoders. Given a minibatch of $N$ paired audio and text samples, denoted by $\{(X_{a,i},X_{t,i})\}_{i=1}^{N}$, the respective feature representations are obtained as

\begin{align}
\hat{X}_a &= f_a(X_a)
\in \mathbb{R}^{N \times V}, \\
\hat{X}_t &= f_t(X_t)
\in \mathbb{R}^{N \times U}.
\end{align}

The audio encoder employs a CNN-based architecture, such as CNN14, operating on log-Mel spectrogram representations, whereas the text encoder utilizes a transformer-based language model such as BERT. The extracted features are subsequently mapped into a common embedding space of dimension $d$ through learnable projection functions:

\begin{align}
E_a &= L_a(\hat{X}_a)
\in \mathbb{R}^{N \times d}, \\
E_t &= L_t(\hat{X}_t)
\in \mathbb{R}^{N \times d}.
\end{align}

Here, $L_a$ and $L_t$ represent the learnable audio and text projection functions, respectively. The projected embeddings are normalized before similarity computation.

\subsubsection{Scaled Similarity Computation}

CLAP computes the pairwise similarity between audio and text embeddings to construct a similarity matrix $C \in \mathbb{R}^{N \times N}$. The scaled similarity is defined as

\begin{equation}
C = \frac{E_t E_a^T}{\tau}
= \exp(t)(E_t E_a^T),
\label{eq:clap_similarity}
\end{equation}

where $\tau$ represents the learnable temperature parameter and $t=\log(1/\tau)$. The inverse temperature controls the sharpness of the similarity distribution. The logit scaling operation is constrained to prevent excessively large similarity values during training.

\subsubsection{Symmetric Contrastive Loss}

CLAP adopts a symmetric cross-entropy objective to align corresponding audio and text representations. The text-to-audio contrastive loss is defined as

\begin{equation}
\ell_{T \rightarrow A}(C)
=
-\frac{1}{N}
\sum_{i=1}^{N}
\log
\frac{\exp(C_{ii})}
{\sum_{j=1}^{N}\exp(C_{ij})}.
\label{eq:clap_t2a}
\end{equation}

Similarly, the audio-to-text contrastive loss is expressed as

\begin{equation}
\ell_{A \rightarrow T}(C)
=
-\frac{1}{N}
\sum_{j=1}^{N}
\log
\frac{\exp(C_{jj})}
{\sum_{i=1}^{N}\exp(C_{ij})}.
\label{eq:clap_a2t}
\end{equation}

The overall CLAP training objective is obtained by averaging the two directional losses:

\begin{equation}
\mathcal{L}_{\mathrm{CLAP}}
=
\frac{1}{2}
\left[
\ell_{T \rightarrow A}(C)
+
\ell_{A \rightarrow T}(C)
\right].
\label{eq:clap_total}
\end{equation}

This objective promotes the alignment of semantically corresponding audio and text embeddings while distinguishing them from non-corresponding samples.

\subsubsection{Zero-Shot Classification}

During zero-shot evaluation, CLAP compares the embedding of an input audio sample against the text embeddings of a predefined set of $K$ candidate class descriptions. Let $C^{*} \in \mathbb{R}^{M \times K}$ denote the resulting scaled similarity matrix for $M$ test audio samples and $K$ candidate text prompts.

For single-label classification, the predicted class probabilities are obtained through the softmax function:

\begin{equation}
\hat{y}_{i,k}
=
\frac{\exp(C^{*}_{i,k})}
{\sum_{j=1}^{K}\exp(C^{*}_{i,j})}.
\label{eq:clap_softmax}
\end{equation}

The predicted class is subsequently determined by selecting the class with the highest probability:

\begin{equation}
\hat{c}_i
=
\underset{k \in \{1,\ldots,K\}}{\operatorname{argmax}}
\ \hat{y}_{i,k}.
\label{eq:clap_prediction}
\end{equation}

For multi-label classification, where multiple acoustic events may coexist within a single audio sample, the sigmoid function is applied independently to each class similarity score:

\begin{equation}
\hat{y}_{i,k}
=
\sigma(C^{*}_{i,k})
=
\frac{1}{1+\exp(-C^{*}_{i,k})}.
\label{eq:clap_sigmoid}
\end{equation}

Unlike softmax, sigmoid does not impose a sum-to-one constraint across class probabilities, allowing multiple classes to be predicted simultaneously. A predefined decision threshold can subsequently be applied to determine the presence or absence of individual acoustic categories.
```

\section{Tables}\label{sec5}

Tables can be inserted via the normal table and tabular environment. To put
footnotes inside tables you should use \verb+\footnotetext[]{...}+ tag.
The footnote appears just below the table itself (refer Tables~\ref{tab1} and \ref{tab2}). 
For the corresponding footnotemark use \verb+\footnotemark[...]+

\begin{table}[h]
\caption{Caption text}\label{tab1}%
\begin{tabular}{@{}llll@{}}
\toprule
Column 1 & Column 2  & Column 3 & Column 4\\
\midrule
row 1    & data 1   & data 2  & data 3  \\
row 2    & data 4   & data 5\footnotemark[1]  & data 6  \\
row 3    & data 7   & data 8  & data 9\footnotemark[2]  \\
\botrule
\end{tabular}
\footnotetext{Source: This is an example of table footnote. This is an example of table footnote.}
\footnotetext[1]{Example for a first table footnote. This is an example of table footnote.}
\footnotetext[2]{Example for a second table footnote. This is an example of table footnote.}
\end{table}

\noindent
The input format for the above table is as follows:

%%=============================================%%
%% For presentation purpose, we have included  %%
%% \bigskip command. Please ignore this.       %%
%%=============================================%%
\bigskip
\begin{verbatim}
\begin{table}[<placement-specifier>]
\caption{<table-caption>}\label{<table-label>}%
\begin{tabular}{@{}llll@{}}
\toprule
Column 1 & Column 2 & Column 3 & Column 4\\
\midrule
row 1 & data 1 & data 2	 & data 3 \\
row 2 & data 4 & data 5\footnotemark[1] & data 6 \\
row 3 & data 7 & data 8	 & data 9\footnotemark[2]\\
\botrule
\end{tabular}
\footnotetext{Source: This is an example of table footnote. 
This is an example of table footnote.}
\footnotetext[1]{Example for a first table footnote.
This is an example of table footnote.}
\footnotetext[2]{Example for a second table footnote. 
This is an example of table footnote.}
\end{table}
\end{verbatim}
\bigskip
%%=============================================%%
%% For presentation purpose, we have included  %%
%% \bigskip command. Please ignore this.       %%
%%=============================================%%

\begin{table}[h]
\caption{Example of a lengthy table which is set to full textwidth}\label{tab2}
\begin{tabular*}{\textwidth}{@{\extracolsep\fill}lcccccc}
\toprule%
& \multicolumn{3}{@{}c@{}}{Element 1\footnotemark[1]} & \multicolumn{3}{@{}c@{}}{Element 2\footnotemark[2]} \\\cmidrule{2-4}\cmidrule{5-7}%
Project & Energy & $\sigma_{calc}$ & $\sigma_{expt}$ & Energy & $\sigma_{calc}$ & $\sigma_{expt}$ \\
\midrule
Element 3  & 990 A & 1168 & $1547\pm12$ & 780 A & 1166 & $1239\pm100$\\
Element 4  & 500 A & 961  & $922\pm10$  & 900 A & 1268 & $1092\pm40$\\
\botrule
\end{tabular*}
\footnotetext{Note: This is an example of table footnote. This is an example of table footnote this is an example of table footnote this is an example of~table footnote this is an example of table footnote.}
\footnotetext[1]{Example for a first table footnote.}
\footnotetext[2]{Example for a second table footnote.}
\end{table}

In case of double column layout, tables which do not fit in single column width should be set to full text width. For this, you need to use \verb+\begin{table*}+ \verb+...+ \verb+\end{table*}+ instead of \verb+\begin{table}+ \verb+...+ \verb+\end{table}+ environment. Lengthy tables which do not fit in textwidth should be set as rotated table. For this, you need to use \verb+\begin{sidewaystable}+ \verb+...+ \verb+\end{sidewaystable}+ instead of \verb+\begin{table*}+ \verb+...+ \verb+\end{table*}+ environment. This environment puts tables rotated to single column width. For tables rotated to double column width, use \verb+\begin{sidewaystable*}+ \verb+...+ \verb+\end{sidewaystable*}+.

\begin{sidewaystable}
\caption{Tables which are too long to fit, should be written using the ``sidewaystable'' environment as shown here}\label{tab3}
\begin{tabular*}{\textheight}{@{\extracolsep\fill}lcccccc}
\toprule%
& \multicolumn{3}{@{}c@{}}{Element 1\footnotemark[1]}& \multicolumn{3}{@{}c@{}}{Element\footnotemark[2]} \\\cmidrule{2-4}\cmidrule{5-7}%
Projectile & Energy	& $\sigma_{calc}$ & $\sigma_{expt}$ & Energy & $\sigma_{calc}$ & $\sigma_{expt}$ \\
\midrule
Element 3 & 990 A & 1168 & $1547\pm12$ & 780 A & 1166 & $1239\pm100$ \\
Element 4 & 500 A & 961  & $922\pm10$  & 900 A & 1268 & $1092\pm40$ \\
Element 5 & 990 A & 1168 & $1547\pm12$ & 780 A & 1166 & $1239\pm100$ \\
Element 6 & 500 A & 961  & $922\pm10$  & 900 A & 1268 & $1092\pm40$ \\
\botrule
\end{tabular*}
\footnotetext{Note: This is an example of table footnote this is an example of table footnote this is an example of table footnote this is an example of~table footnote this is an example of table footnote.}
\footnotetext[1]{This is an example of table footnote.}
\end{sidewaystable}

\section{Figures}\label{sec6}

As per the \LaTeX\ standards you need to use eps images for \LaTeX\ compilation and \verb+pdf/jpg/png+ images for \verb+PDFLaTeX+ compilation. This is one of the major difference between \LaTeX\ and \verb+PDFLaTeX+. Each image should be from a single input .eps/vector image file. Avoid using subfigures. The command for inserting images for \LaTeX\ and \verb+PDFLaTeX+ can be generalized. The package used to insert images in \verb+LaTeX/PDFLaTeX+ is the graphicx package. Figures can be inserted via the normal figure environment as shown in the below example:

%%=============================================%%
%% For presentation purpose, we have included  %%
%% \bigskip command. Please ignore this.       %%
%%=============================================%%
\bigskip
\begin{verbatim}
\begin{figure}[<placement-specifier>]
\centering
\includegraphics{<eps-file>}
\caption{<figure-caption>}\label{<figure-label>}
\end{figure}
\end{verbatim}
\bigskip
%%=============================================%%
%% For presentation purpose, we have included  %%
%% \bigskip command. Please ignore this.       %%
%%=============================================%%

\begin{figure}[h]
\centering
\includegraphics[width=0.9\textwidth]{fig.eps}
\caption{This is a widefig. This is an example of long caption this is an example of long caption  this is an example of long caption this is an example of long caption}\label{fig1}
\end{figure}

In case of double column layout, the above format puts figure captions/images to single column width. To get spanned images, we need to provide \verb+\begin{figure*}+ \verb+...+ \verb+\end{figure*}+.

For sample purpose, we have included the width of images in the optional argument of \verb+\includegraphics+ tag. Please ignore this. 

\section{Algorithms, Program codes and Listings}\label{sec7}

Packages \verb+algorithm+, \verb+algorithmicx+ and \verb+algpseudocode+ are used for setting algorithms in \LaTeX\ using the format:

%%=============================================%%
%% For presentation purpose, we have included  %%
%% \bigskip command. Please ignore this.       %%
%%=============================================%%
\bigskip
\begin{verbatim}
\begin{algorithm}
\caption{<alg-caption>}\label{<alg-label>}
\begin{algorithmic}[1]
. . .
\end{algorithmic}
\end{algorithm}
\end{verbatim}
\bigskip
%%=============================================%%
%% For presentation purpose, we have included  %%
%% \bigskip command. Please ignore this.       %%
%%=============================================%%

You may refer above listed package documentations for more details before setting \verb+algorithm+ environment. For program codes, the ``verbatim'' package is required and the command to be used is \verb+\begin{verbatim}+ \verb+...+ \verb+\end{verbatim}+. 

Similarly, for \verb+listings+, use the \verb+listings+ package. \verb+\begin{lstlisting}+ \verb+...+ \verb+\end{lstlisting}+ is used to set environments similar to \verb+verbatim+ environment. Refer to the \verb+lstlisting+ package documentation for more details.

A fast exponentiation procedure:

\lstset{texcl=true,basicstyle=\small\sf,commentstyle=\small\rm,mathescape=true,escapeinside={(*}{*)}}
\begin{lstlisting}
begin
  for $i:=1$ to $10$ step $1$ do
      expt($2,i$);  
      newline() od                (*\textrm{Comments will be set flush to the right margin}*)
where
proc expt($x,n$) $\equiv$
  $z:=1$;
  do if $n=0$ then exit fi;
     do if odd($n$) then exit fi;                 
        comment: (*\textrm{This is a comment statement;}*)
        $n:=n/2$; $x:=x*x$ od;
     { $n>0$ };
     $n:=n-1$; $z:=z*x$ od;
  print($z$). 
end
\end{lstlisting}

\begin{algorithm}
\caption{Calculate $y = x^n$}\label{algo1}
\begin{algorithmic}[1]
\Require $n \geq 0 \vee x \neq 0$
\Ensure $y = x^n$ 
\State $y \Leftarrow 1$
\If{$n < 0$}\label{algln2}
        \State $X \Leftarrow 1 / x$
        \State $N \Leftarrow -n$
\Else
        \State $X \Leftarrow x$
        \State $N \Leftarrow n$
\EndIf
\While{$N \neq 0$}
        \If{$N$ is even}
            \State $X \Leftarrow X \times X$
            \State $N \Leftarrow N / 2$
        \Else[$N$ is odd]
            \State $y \Leftarrow y \times X$
            \State $N \Leftarrow N - 1$
        \EndIf
\EndWhile
\end{algorithmic}
\end{algorithm}

%%=============================================%%
%% For presentation purpose, we have included  %%
%% \bigskip command. Please ignore this.       %%
%%=============================================%%
\bigskip
\begin{minipage}{\hsize}%
\lstset{frame=single,framexleftmargin=-1pt,framexrightmargin=-17pt,framesep=12pt,linewidth=0.98\textwidth,language=pascal}% Set your language (you can change the language for each code-block optionally)
%%% Start your code-block
\begin{lstlisting}
for i:=maxint to 0 do
begin
{ do nothing }
end;
Write('Case insensitive ');
Write('Pascal keywords.');
\end{lstlisting}
\end{minipage}

\section{Cross referencing}\label{sec8}

Environments such as figure, table, equation and align can have a label
declared via the \verb+\label{#label}+ command. For figures and table
environments use the \verb+\label{}+ command inside or just
below the \verb+\caption{}+ command. You can then use the
\verb+\ref{#label}+ command to cross-reference them. As an example, consider
the label declared for Figure~\ref{fig1} which is
\verb+\label{fig1}+. To cross-reference it, use the command 
\verb+Figure \ref{fig1}+, for which it comes up as
``Figure~\ref{fig1}''. 

To reference line numbers in an algorithm, consider the label declared for the line number 2 of Algorithm~\ref{algo1} is \verb+\label{algln2}+. To cross-reference it, use the command \verb+\ref{algln2}+ for which it comes up as line~\ref{algln2} of Algorithm~\ref{algo1}.

\subsection{Details on reference citations}\label{subsec7}

Standard \LaTeX\ permits only numerical citations. To support both numerical and author-year citations this template uses \verb+natbib+ \LaTeX\ package. For style guidance please refer to the template user manual.

Here is an example for \verb+\cite{...}+: \cite{bib1}. Another example for \verb+\citep{...}+: \citep{bib2}. For author-year citation mode, \verb+\cite{...}+ prints Jones et al. (1990) and \verb+\citep{...}+ prints (Jones et al., 1990).

All cited bib entries are printed at the end of this article: \cite{bib3}, \cite{bib4}, \cite{bib5}, \cite{bib6}, \cite{bib7}, \cite{bib8}, \cite{bib9}, \cite{bib10}, \cite{bib11}, \cite{bib12} and \cite{bib13}.

\section{Examples for theorem like environments}\label{sec10}

For theorem like environments, we require \verb+amsthm+ package. There are three types of predefined theorem styles exists---\verb+thmstyleone+, \verb+thmstyletwo+ and \verb+thmstylethree+ 

%%=============================================%%
%% For presentation purpose, we have included  %%
%% \bigskip command. Please ignore this.       %%
%%=============================================%%
\bigskip
\begin{tabular}{|l|p{19pc}|}
\hline
\verb+thmstyleone+ & Numbered, theorem head in bold font and theorem text in italic style \\\hline
\verb+thmstyletwo+ & Numbered, theorem head in roman font and theorem text in italic style \\\hline
\verb+thmstylethree+ & Numbered, theorem head in bold font and theorem text in roman style \\\hline
\end{tabular}
\bigskip
%%=============================================%%
%% For presentation purpose, we have included  %%
%% \bigskip command. Please ignore this.       %%
%%=============================================%%

For mathematics journals, theorem styles can be included as shown in the following examples:

\begin{theorem}[Theorem subhead]\label{thm1}
Example theorem text. Example theorem text. Example theorem text. Example theorem text. Example theorem text. 
Example theorem text. Example theorem text. Example theorem text. Example theorem text. Example theorem text. 
Example theorem text. 
\end{theorem}

Sample body text. Sample body text. Sample body text. Sample body text. Sample body text. Sample body text. Sample body text. Sample body text.

\begin{proposition}
Example proposition text. Example proposition text. Example proposition text. Example proposition text. Example proposition text. 
Example proposition text. Example proposition text. Example proposition text. Example proposition text. Example proposition text. 
\end{proposition}

Sample body text. Sample body text. Sample body text. Sample body text. Sample body text. Sample body text. Sample body text. Sample body text.

\begin{example}
Phasellus adipiscing semper elit. Proin fermentum massa
ac quam. Sed diam turpis, molestie vitae, placerat a, molestie nec, leo. Maecenas lacinia. Nam ipsum ligula, eleifend
at, accumsan nec, suscipit a, ipsum. Morbi blandit ligula feugiat magna. Nunc eleifend consequat lorem. 
\end{example}

Sample body text. Sample body text. Sample body text. Sample body text. Sample body text. Sample body text. Sample body text. Sample body text.

\begin{remark}
Phasellus adipiscing semper elit. Proin fermentum massa
ac quam. Sed diam turpis, molestie vitae, placerat a, molestie nec, leo. Maecenas lacinia. Nam ipsum ligula, eleifend
at, accumsan nec, suscipit a, ipsum. Morbi blandit ligula feugiat magna. Nunc eleifend consequat lorem. 
\end{remark}

Sample body text. Sample body text. Sample body text. Sample body text. Sample body text. Sample body text. Sample body text. Sample body text.

\begin{definition}[Definition sub head]
Example definition text. Example definition text. Example definition text. Example definition text. Example definition text. Example definition text. Example definition text. Example definition text. 
\end{definition}

Additionally a predefined ``proof'' environment is available: \verb+\begin{proof}+ \verb+...+ \verb+\end{proof}+. This prints a ``Proof'' head in italic font style and the ``body text'' in roman font style with an open square at the end of each proof environment. 

\begin{proof}
Example for proof text. Example for proof text. Example for proof text. Example for proof text. Example for proof text. Example for proof text. Example for proof text. Example for proof text. Example for proof text. Example for proof text. 
\end{proof}

Sample body text. Sample body text. Sample body text. Sample body text. Sample body text. Sample body text. Sample body text. Sample body text.

\begin{proof}[Proof of Theorem~{\upshape\ref{thm1}}]
Example for proof text. Example for proof text. Example for proof text. Example for proof text. Example for proof text. Example for proof text. Example for proof text. Example for proof text. Example for proof text. Example for proof text. 
\end{proof}

\noindent
For a quote environment, use \verb+\begin{quote}...\end{quote}+
\begin{quote}
Quoted text example. Aliquam porttitor quam a lacus. Praesent vel arcu ut tortor cursus volutpat. In vitae pede quis diam bibendum placerat. Fusce elementum
convallis neque. Sed dolor orci, scelerisque ac, dapibus nec, ultricies ut, mi. Duis nec dui quis leo sagittis commodo.
\end{quote}

Sample body text. Sample body text. Sample body text. Sample body text. Sample body text (refer Figure~\ref{fig1}). Sample body text. Sample body text. Sample body text (refer Table~\ref{tab3}). 

\section{Methods}\label{sec11}

Topical subheadings are allowed. Authors must ensure that their Methods section includes adequate experimental and characterization data necessary for others in the field to reproduce their work. Authors are encouraged to include RIIDs where appropriate. 

\textbf{Ethical approval declarations} (only required where applicable) Any article reporting experiment/s carried out on (i)~live vertebrate (or higher invertebrates), (ii)~humans or (iii)~human samples must include an unambiguous statement within the methods section that meets the following requirements: 

\begin{enumerate}[1.]
\item Approval: a statement which confirms that all experimental protocols were approved by a named institutional and/or licensing committee. Please identify the approving body in the methods section

\item Accordance: a statement explicitly saying that the methods were carried out in accordance with the relevant guidelines and regulations

\item Informed consent (for experiments involving humans or human tissue samples): include a statement confirming that informed consent was obtained from all participants and/or their legal guardian/s
\end{enumerate}

If your manuscript includes potentially identifying patient/participant information, or if it describes human transplantation research, or if it reports results of a clinical trial then  additional information will be required. Please visit (\url{https://www.nature.com/nature-research/editorial-policies}) for Nature Portfolio journals, (\url{https://www.springer.com/gp/authors-editors/journal-author/journal-author-helpdesk/publishing-ethics/14214}) for Springer Nature journals, or (\url{https://www.biomedcentral.com/getpublished/editorial-policies\#ethics+and+consent}) for BMC.

\section{Discussion}\label{sec12}

Discussions should be brief and focused. In some disciplines use of Discussion or `Conclusion' is interchangeable. It is not mandatory to use both. Some journals prefer a section `Results and Discussion' followed by a section `Conclusion'. Please refer to Journal-level guidance for any specific requirements. 

\section{Conclusion}\label{sec13}

Conclusions may be used to restate your hypothesis or research question, restate your major findings, explain the relevance and the added value of your work, highlight any limitations of your study, describe future directions for research and recommendations. 

In some disciplines use of Discussion or 'Conclusion' is interchangeable. It is not mandatory to use both. Please refer to Journal-level guidance for any specific requirements. 

\backmatter

\bmhead{Supplementary information}

If your article has accompanying supplementary file/s please state so here. 

Authors reporting data from electrophoretic gels and blots should supply the full unprocessed scans for key as part of their Supplementary information. This may be requested by the editorial team/s if it is missing.

Please refer to Journal-level guidance for any specific requirements.

\bmhead{Acknowledgements}

Acknowledgements are not compulsory. Where included they should be brief. Grant or contribution numbers may be acknowledged.

Please refer to Journal-level guidance for any specific requirements.

\section*{Declarations}

Some journals require declarations to be submitted in a standardised format. Please check the Instructions for Authors of the journal to which you are submitting to see if you need to complete this section. If yes, your manuscript must contain the following sections under the heading `Declarations':

\begin{itemize}
\item Funding
\item Conflict of interest/Competing interests (check journal-specific guidelines for which heading to use)
\item Ethics approval and consent to participate
\item Consent for publication
\item Data availability 
\item Materials availability
\item Code availability 
\item Author contribution
\end{itemize}

\noindent
If any of the sections are not relevant to your manuscript, please include the heading and write `Not applicable' for that section. 

%%===================================================%%
%% For presentation purpose, we have included        %%
%% \bigskip command. Please ignore this.             %%
%%===================================================%%
\bigskip
\begin{flushleft}%
Editorial Policies for:

\bigskip\noindent
Springer journals and proceedings: \url{https://www.springer.com/gp/editorial-policies}

\bigskip\noindent
Nature Portfolio journals: \url{https://www.nature.com/nature-research/editorial-policies}

\bigskip\noindent
\textit{Scientific Reports}: \url{https://www.nature.com/srep/journal-policies/editorial-policies}

\bigskip\noindent
BMC journals: \url{https://www.biomedcentral.com/getpublished/editorial-policies}
\end{flushleft}

\begin{appendices}

\section{Section title of first appendix}\label{secA1}

An appendix contains supplementary information that is not an essential part of the text itself but which may be helpful in providing a more comprehensive understanding of the research problem or it is information that is too cumbersome to be included in the body of the paper.

%%=============================================%%
%% For submissions to Nature Portfolio Journals %%
%% please use the heading ``Extended Data''.   %%
%%=============================================%%

%%=============================================================%%
%% Sample for another appendix section			       %%
%%=============================================================%%

%% \section{Example of another appendix section}\label{secA2}%
%% Appendices may be used for helpful, supporting or essential material that would otherwise 
%% clutter, break up or be distracting to the text. Appendices can consist of sections, figures, 
%% tables and equations etc.

\end{appendices}

%%===========================================================================================%%
%% If you are submitting to one of the Nature Portfolio journals, using the eJP submission   %%
%% system, please include the references within the manuscript file itself. You may do this  %%
%% by copying the reference list from your .bbl file, paste it into the main manuscript .tex %%
%% file, and delete the associated \verb+\bibliography+ commands.                            %%
%%===========================================================================================%%

\bibliography{sn-bibliography}% common bib file
%% if required, the content of .bbl file can be included here once bbl is generated
%%\input sn-article.bbl

\end{document}
